% SPDX-FileCopyrightText: 2026 Sascha Brawer <sascha@brawer.ch>
% SPDX-License-Identifier: CC-BY-4.0
%
% Build:  make        (LuaLaTeX, TeX Live 2025+; runs twice, and pins the PDF
%                       timestamps via SOURCE_DATE_EPOCH for a reproducible build)
% Verify: make check   (verapdf: PDF/UA-2 + PDF/A-4f)
% Output: defensive-publication.pdf -- a tagged PDF 2.0.
%
% This .tex is the single source. The publication date is set in ONE place --
% \pubdate below -- and must match PUB_DATE in the Makefile.

\DocumentMetadata{
  lang        = en-US,
  pdfversion  = 2.0,
  pdfstandard = ua-2,
  pdfstandard = A-4,   % -> becomes PDF/A-4f, since luamml embeds MathML files
  testphase   = {phase-III, math, graphic},
}

\documentclass[10pt,twocolumn]{article}

%--- Fonts: Source Serif 4 (text) + Libertinus Math + Source Code Pro (mono),
%    the two Adobe Source faces being a designed pair. In TeX Live the family
%    is still packaged under its former name, "Source Serif Pro".
\usepackage{fontspec}
\setmainfont{Source Serif Pro}[Ligatures=TeX,Numbers=Lining]
\setmonofont{Source Code Pro}[Scale=MatchLowercase]
\usepackage{amsmath}
\usepackage{unicode-math}
\setmathfont{Libertinus Math}
\usepackage[american]{babel}          % US hyphenation patterns
\usepackage[autostyle=true]{csquotes} % \enquote -> American quotation marks
\usepackage{luamml} % emit MathML for the equations (math accessibility)

%--- Page geometry: A4, disclosure-style two-column block ------------------
\usepackage[a4paper,margin=19mm,columnsep=5mm]{geometry}
\usepackage{microtype}
\usepackage{graphicx} % for the CC-BY licence badge (vector, bundled with doclicense)
\setlength{\emergencystretch}{2em}

% Tight list spacing. The phase-III tag pipeline reimplements lists and bypasses
% enumitem, but still reads the classic \@list.. glue parameters, so set those.
\makeatletter
\renewcommand\@listi{%
  \leftmargin\leftmargini
  \topsep    3\p@ \@plus\p@ \@minus\p@
  \parsep    \z@
  \itemsep   2\p@ \@plus\p@}
\let\@listI\@listi
\renewcommand\@listii{%
  \leftmargin\leftmarginii \labelwidth\leftmarginii \advance\labelwidth-\labelsep
  \topsep 2\p@ \@plus\p@ \@minus\p@ \parsep \z@ \itemsep \z@}
\makeatother

%--- Headings -------------------------------------------------------------
\usepackage{titlesec}
\titleformat{\section}{\normalfont\large\bfseries}{\thesection}{0.6em}{}
\titleformat{\subsection}{\normalfont\normalsize\bfseries}{\thesubsection}{0.6em}{}
\titleformat{\subsubsection}{\normalfont\normalsize\itshape}{\thesubsubsection}{0.6em}{}
\titlespacing*{\section}{0pt}{1.6ex plus .3ex minus .2ex}{0.7ex}
\titlespacing*{\subsection}{0pt}{1.3ex plus .3ex minus .2ex}{0.5ex}
\titlespacing*{\subsubsection}{0pt}{1.1ex plus .2ex minus .2ex}{0.4ex}

%--- Publication date: single source of truth (keep Makefile PUB_DATE in sync)
\newcommand{\pubdate}{September 2, 2026} % visible masthead form
\newcommand{\pubdateiso}{2026-09-02}     % metadata form (comma-free: the XMP
                                         % writer splits pdfsubject on commas)

%--- Links and document metadata (feeds the XMP packet, hence indexing) -----
\usepackage{hyperref}
\hypersetup{
  % {,} protects the comma: the XMP writer otherwise splits the title into a list
  pdftitle    = {Method for Memory-Bounded Construction of a Globally Complete{,} High-Zoom-Level Cloud-Optimized GeoTIFF from Tile-Access Logs},
  pdfauthor   = {Sascha Brawer},
  pdfsubject  = {Defensive publication (prior-art disclosure). Published on Technical Disclosure Commons on \pubdateiso.},
  pdfkeywords = {defensive publication, prior art, technical disclosure, Cloud-Optimized GeoTIFF, COG, TIFF, space-filling curve, Morton code, Z-order curve, external sort, k-way merge, streaming aggregation, windowed median, overview pyramid, mipmap, tile deduplication, shared file offsets, web map tiles, tile-access logs, log analysis, OpenStreetMap, constant memory, Web Mercator},
  colorlinks  = true,
  linkcolor   = black,
  citecolor   = black,
  urlcolor    = {blue!45!black},
  pdflicenseurl = {https://creativecommons.org/licenses/by/4.0/},
}

\newcommand{\code}[1]{\texttt{#1}}

% Unnumbered first-page footnote (conference-style copyright/venue note).
% \footnotetext places no mark in the running text; the empty \thefootnote
% keeps a stray "0" from appearing at the foot.
\newcommand\blfootnote[1]{%
  \begingroup\renewcommand\thefootnote{}\footnotetext{#1}\endgroup
}

% Drop the blank 1.8em mark box that the phase-III footnote code puts at the
% start of every footnote -- with an empty mark it is just an indent, and it
% was pushing the licence badge away from the column edge.
\setlength\footnotemargin{0pt}

% First-page licence footnote: the CC badge sits flush against the column's
% left edge (aligned with the flush lines of the note); the first two lines
% of the note are hung to just past the badge's right edge -- no other indent
% -- and the rest of the note reclaims the full column width.
%
% NB: pre-compute the indent into a register and use bare \hangindent\ccind /
% \hangafter-2 -- writing \hangindent=\dimexpr...\relax on the same line as
% \hangafter mis-scans and only line 1 ends up hung.
\newsavebox\ccbadgebox
\newlength\ccgap \setlength\ccgap{0.6em} % badge-to-text gap only
\newlength\ccind
\newcommand\cclicensefootnote[1]{%
  \blfootnote{%
    \footnotesize
    \sbox\ccbadgebox{\includegraphics[height=1.55\baselineskip,%
      alt={Creative Commons Attribution 4.0 International (CC BY 4.0) license badge}]%
      {doclicense-CC-by-88x31}}%
    \setlength\ccind{\dimexpr\wd\ccbadgebox+\ccgap\relax}%
    \hangindent\ccind \hangafter-2
    \noindent
    \hspace*{-\ccind}%
    % bottom edge of the badge on the baseline of the note's second line
    \raisebox{-\baselineskip}[0pt][0pt]{\usebox\ccbadgebox}%
    \hspace*{\dimexpr\ccind-\wd\ccbadgebox\relax}%
    \ignorespaces #1\par
  }%
}

\setlength{\parskip}{0pt}
\setlength{\parindent}{1.2em}
\frenchspacing

\begin{document}
\raggedbottom % after the class re-asserts \flushbottom for twocolumn; no vertical stretch

%========================================================================
% Full-width masthead
%========================================================================
\twocolumn[{%
  \begin{center}
    {\LARGE\bfseries Method for Memory-Bounded Construction of a\\[3pt]
     Globally Complete, High-Zoom-Level Cloud-Optimized\\[3pt]
     GeoTIFF from Tile-Access Logs\par}
    \vspace{1em}
    {\large Sascha Brawer\,\textperiodcentered\,\href{mailto:sascha@brawer.ch}{sascha@brawer.ch}\,\textperiodcentered\,Bern, Switzerland\par}
    \vspace{0.35em}
    {\small Published on Technical Disclosure Commons
     (\url{https://www.tdcommons.org/})\,\textperiodcentered\,\pubdate\par}
  \end{center}
  \vspace{0.9em}
  \hrule
  \vspace{0.9em}
  \begingroup\small
  \noindent\textbf{Abstract.}
  A memory-bounded method is disclosed for constructing a globally complete,
  high-zoom-level Cloud-Optimized GeoTIFF --- a planet-wide raster at web-map
  zoom level~18, on the order of $10^{11}$ cells --- from a rolling window of
  tile-access logs, using working memory that is essentially constant in the
  number of cells. The output is a standards-compliant COG with a full
  reduced-resolution overview pyramid, derived by robustly aggregating the
  per-tile time series; it is several hundred times smaller than a na\"ive
  dense encoding. The method combines a level-embedding space-filling tile key,
  a two-stage disk-backed ordering (per-period external sort followed by a
  cross-period streaming merge), a bounded root-to-leaf raster working set that
  follows from consuming the ordered stream in tree pre-order, single-pass
  inline overview generation, per-cell robust aggregation computed without
  materializing the time series, and three complementary output-size reductions
  (magnitude quantization, tile deduplication by shared file offsets, and a
  monotone value transform). A worked configuration builds a global zoom-18
  raster in under one gigabyte of resident memory and emits a roughly 600~MB
  file where a dense encoding would require approximately 275~GB.\par
  \medskip
  \noindent\textbf{Keywords.} Defensive publication, prior art, Cloud-Optimized
  GeoTIFF, space-filling curve, Morton / Z-order code, external sort, streaming
  $k$-way merge, windowed median, single-pass overview pyramid, tile
  deduplication by shared file offsets, tile-access logs, log analysis, web map
  tiles, constant-memory raster construction, OpenStreetMap.\par
  \medskip
  \noindent\textbf{Purpose.} This document is published to establish prior art.
  The techniques described here, and obvious variations of them, are disclosed
  for free public use; this publication is intended to prevent anyone from
  obtaining patent rights that would restrict that use.\par
  \medskip
  \noindent\textbf{Reference implementation.}
  \url{https://github.com/brawer/osmviews} (\code{cmd/osmviews-builder}),
  released under the MIT license; see~\S6.\par
  \endgroup
  \vspace{0.9em}
  \hrule
  \vspace{1.4em}
}]

% Conference-style first-page footnote: licence, AI declaration.
\cclicensefootnote{%
  The text of this document is licensed under a Creative Commons Attribution
  4.0 International License (CC BY 4.0),
  \url{https://creativecommons.org/licenses/by/4.0/}; its source is in the same
  repository as the reference implementation (\S6), under
  \code{docs/defensive-publication}. The reference implementation is licensed
  separately, under the MIT license.
  The method disclosed here is the author's own, developed without AI-assisted
  research; the reference implementation includes minor AI-assisted code
  changes, marked in the repository's commit history. This document was
  drafted, copy-edited, and typeset by an AI assistant (Anthropic's Claude)
  from the author's technical input, under the author's direction and review.%
}

%========================================================================
\section{Problem}
%========================================================================

Consider a data set of \emph{tile-access counts}: for a web map tiled in the
usual \code{zoom/x/y} scheme, a stream of records stating how many times each
tile was requested in a given period (e.g.\ one day). Such logs are published,
for example, by the OpenStreetMap Foundation.

The goal is to convert a rolling window of such logs (e.g.\ the most recent 52
weekly periods) into a geospatial raster in which each cell holds an
aggregated, area-normalized \enquote{view density} for the corresponding
location, at a spatial resolution equivalent to zoom level~18 --- roughly
150~m per cell at the equator --- with coverage of the entire globe, and
delivered as a Cloud-Optimized GeoTIFF (COG) so that clients can read arbitrary
sub-regions and zoom levels over HTTP range requests.

The na\"ive approach --- allocate the full dense zoom-18 array, accumulate into
it, then write it --- is infeasible on commodity hardware. A global zoom-18
raster has $2^{18} \times 2^{18} \approx 6.9 \times 10^{10}$ cells; at four
bytes per cell that is approximately 275~GB, before overviews. The aggregation
also requires holding, per cell, enough of the time series to compute a robust
statistic such as a median. Neither fits in memory, and materializing either on
disk in dense form is wasteful because the overwhelming majority of the planet
(oceans, deserts, ice) has little or no data.

The techniques below build such a raster in a single streaming pass, with peak
memory that scales with the \emph{depth} of the tile tree (a small constant,
e.g.\ 11 or 19) rather than its \emph{breadth} (the cell count), and with an
output file whose size tracks the information content of the data rather than
the dimensions of the grid.

%========================================================================
\section{Overview of the approach}
%========================================================================

\begin{enumerate}
\item Encode each \code{zoom/x/y} tile as an integer \emph{tile key} whose
  numeric order is a depth-first pre-order traversal of the tile quad-tree
  (\S3.1).
\item For each time period, sort that period's records by tile key using an
  external (disk-backed) sort, and coalesce equal keys; store the sorted,
  coalesced period as a compressed file (\S3.2).
\item Merge the sorted period files with a streaming $k$-way merge, so that a
  single globally key-ordered record stream is produced, in which every
  period's records for a given tile arrive consecutively (\S3.2).
\item Consume that stream. For each tile, collect its per-period values and
  reduce them to one aggregated value with a robust statistic that tolerates
  missing periods, computed directly from the ordered run (\S3.3).
\item Paint aggregated values into $256 \times 256$ raster tiles held in a
  tree. Because the stream is in pre-order, only the chain of raster tiles from
  the root to the tile currently being filled is ever live; completed raster
  tiles are evicted immediately (\S3.4).
\item On eviction, subsample the completed raster tile into its parent,
  building the overview pyramid inline during the same pass (\S3.5).
\item When writing each raster tile, apply three complementary size reductions:
  quantize low-magnitude tiles to a common value; store each distinct tile body
  once and make other tiles reference the same file bytes; and store a monotone
  transform of the value (\S3.6).
\item Assemble the COG: image file directories first, tile bodies staged in a
  temporary file, offsets back-patched (\S3.7).
\end{enumerate}

Taken individually, several ingredients are well known --- Morton / Z-order
curves, external sorting, $k$-way merges, mipmap generation. The disclosure is
the specific integration that yields constant-memory, single-pass construction
of a \emph{globally complete, high-zoom, overview-complete COG from a time
series}, together with the size-reduction techniques of \S3.6, in particular
tile deduplication by shared file offsets (\S3.6.2).

%========================================================================
\section{Detailed description}
%========================================================================

\subsection{Level-embedding space-filling tile key}

Each tile $(z, x, y)$ is mapped to a fixed-width unsigned integer key as
follows. Reserve the least-significant $b$ bits for the zoom level $z$
($b = 5$ holds any $z \le 31$). Above those, interleave the coordinate bits
two per level, most-significant first: bits $z - 1$ of $x$ and $y$ (the
root-quadrant choice) fill the top key-bit pair, bits $z - 2$ the next, and so
on down to bits $0$ in the lowest pair. Unused low bits are zero. The
coordinate field is then $2z$ bits wide, so a 64-bit key stays collision-free
for $z \le 29$; the reference implementation accepts $z \le 24$.

This is an interleaving (Z-order / Morton) code with two properties that the
method depends on:

\begin{itemize}
\item \textbf{Numeric order equals pre-order traversal.} For two keys, the
  comparison proceeds level by level from the root. A tile's key is numerically
  smaller than the keys of all tiles strictly inside it, and an entire sub-tree
  occupies one contiguous key interval. Sorting keys ascending therefore yields
  a depth-first, pre-order enumeration of the quad-tree.
\item \textbf{Cheap tree operations.} \enquote{Is tile $A$ an ancestor of tile
  $B$}, \enquote{truncate $B$ to level $k$}, and \enquote{successor of $A$ in
  pre-order at maximum level $m$} are all constant-time bit manipulations of
  the key.
\end{itemize}

Logs contain tiles at many zoom levels, so the record stream is
\emph{mixed-level}; embedding the level in the key --- rather than interleaving
$x$ and $y$ alone at a fixed level --- is what makes such a stream sortable
into a single valid traversal.

\emph{Variations.} Any bijective encoding with the pre-order property works:
the level bits may occupy the high end instead of the low end; a Hilbert-curve
ordering may be substituted where its locality is preferred and the ancestor
test is adjusted; trees of branching factor other than four (e.g.\ binary
\enquote{tile pyramids}, octrees for volumetric data) generalize directly by
changing the bits-per-level.

\subsection{Two-stage ordering: per-period external sort, cross-period
streaming merge}

\textbf{Per period.} The records of one period are parsed to (tile key, count)
pairs and fed to an external (disk-backed) sort on that key. The sort spills
runs to temporary files and merges them, so peak memory is bounded by its
configured buffer regardless of the number of distinct tiles. The sorted
output is streamed once more to coalesce equal keys (a tile requested on
several days of the period becomes one record with the summed count) and
written to a compressed file. These per-period files are
cached (locally and, optionally, in object storage) and reused on subsequent
runs, so that a periodic rebuild only sorts the newest period(s).

\textbf{Across periods.} To aggregate the window, the $N$ sorted period
files are combined by a streaming $k$-way merge: maintain a min-heap of
$N$ cursors ordered by tile key, then repeatedly emit the smallest and
advance its cursor. The result is one record stream, globally ordered by tile
key, in which
\textbf{all $N$ periods' records for a given tile arrive consecutively},
followed by all records for the next tile, in quad-tree pre-order. Working
memory for the merge is the heap plus one buffered record per period ---
constant in the cell count and linear only in the (small, fixed) window
size~$N$.

An optional per-period scalar weight may be applied to counts as they are read
from each cursor (for example, to extrapolate a period with missing days to a
full period by scaling by the ratio of expected to observed days); this does
not affect ordering.

In an implementation, the merge and the paint stage run as a two-stage
pipeline: a single \textbf{producer} runs the merge, decompressing the period
files as it reads them, while a single \textbf{consumer} (\S3.3--\S3.7) paints
the raster. The two are coupled by a bounded buffer, so decompression and
painting overlap and the buffer adds only a constant to working memory; each
stage is single-threaded.

\subsection{Per-cell aggregation from the ordered stream}

The consumer reads the merged stream in order and accumulates the run of
records that share a tile key into a small buffer (at most $N$ values).
When the key changes, it reduces the buffer to one value.

The \textbf{windowed median} is a robust reduction that tolerates missing
periods and is computable directly from the ordered run: treat periods with no
record for this tile as zeros, conceptually prepended to the sorted list of
present values. Let $c$ be the number of present values (periods that have
a record for this tile). The median is then the element at index
$\lfloor N/2 \rfloor$ of the length-$N$ conceptual list, i.e.\ the present value
at index $\lfloor N/2 \rfloor - (N - c)$, or zero if that index is negative
(the tile appears in fewer than half the periods). Because
the $k$-way merge breaks tile-key ties by ascending count, the buffered present
values already arrive sorted, so no per-cell sort is needed. The median
suppresses transient spikes, seasonality, and the occasional noisy weighted
period.

The reduced value is then divided by the tile's true surface area (which, in
Web Mercator, varies with latitude) to obtain an area-normalized density.

\emph{Variations.} Any order statistic (a different quantile, a trimmed mean),
or a decay-weighted mean, may be substituted; the point is that the statistic is
obtained from the length-$c$ ascending run plus the count of absent periods
($N - c$), without ever holding all $N$ periods for all cells.

\subsection{Bounded raster working set via pre-order consumption}

The output raster is produced in fixed-size \textbf{raster tiles} of
$T \times T$ cells (e.g.\ $T = 256$). The full-resolution level is a grid of
$2^{L} \times 2^{L}$ raster tiles, where $L = \text{zoom} - \log_2(T)$
(e.g.\ $L = 10$ for zoom~18, $T = 256$). A raster tile is a $T \times T$ array
of the cell type plus a pointer to its parent raster tile.

The consumer maintains only a \textbf{path of raster tiles from the root to the
one currently being filled} --- at most $L + 1$ of them. When a record arrives
whose target raster tile is not contained in the current one, the pre-order
property guarantees that no later record can fall inside the current raster
tile or inside any ancestor that does not contain the new target. Therefore
those raster tiles are \emph{final}. The consumer:

\begin{enumerate}
\item evicts each finalized raster tile (compress its body, append it to the
  output staging file, record its offset and length; see \S3.5 and \S3.6);
\item walks down toward the new target, allocating the chain of intermediate
  raster tiles, and for every sibling sub-tree that is skipped entirely, emits
  a single uniform raster tile (\S3.6.1) carrying a fill value (zero, or a
  nearest-neighbor carry-forward);
\item allocates the new leaf raster tile.
\end{enumerate}

Consequently, the peak live raster memory is
$(L + 1) \cdot T^{2} \cdot \operatorname{sizeof}(\text{cell})$ --- for
$L = 10$, $T = 256$, four-byte cells, about 2.8~MiB --- \textbf{independent of
how many cells the raster has}. Peak process memory in a worked implementation
is about 800~MiB, dominated by language-runtime overhead, compression buffers,
and the $N$ open period readers (a decompressor and one buffered record each),
not by the raster.

Regions never mentioned by any record are emitted as uniform raster tiles
during a final sweep, so the raster has no holes: it covers the whole grid.

\subsection{Single-pass inline overview generation}

When a leaf or interior raster tile is finalized and about to be evicted, it is
first subsampled $2 \times 2 \to 1$ into the appropriate quadrant of its parent
raster tile (e.g.\ by max-pooling, which preserves visible hotspots; average-
or min-pooling are alternatives). Because eviction happens in pre-order and a
parent is only finalized after all four of its children, the entire overview
pyramid (levels $L - 1$ down to $0$) is produced \textbf{within the same single
pass} over the input, with no separate downsampling stage and no re-reading of
the full-resolution data.

\subsection{Output-size reduction}

Three techniques act together so that the file size tracks the data's
information content rather than the grid dimensions.

\subsubsection{Magnitude quantization}

Before compressing a raster tile, test whether every cell, \emph{rounded to the
storage precision} (e.g.\ to the nearest integer density), equals the same
value. Because the great majority of the planet has a density far below one
unit, huge contiguous areas round to a single value (typically zero). A raster
tile that is uniform after rounding is handled as a uniform tile keyed by that
rounded value. This both removes the need to compress most tiles and feeds the
deduplication of \S3.6.2.

\subsubsection{Tile deduplication by shared file offsets}

Maintain, per pyramid level, a map from \emph{tile body content} to the file
offset and length at which a tile with that content was first written. In the
reference implementation the key is the shared rounded value of a uniform tile
(\S3.6.1); hashing the compressed body would extend the same mechanism to
non-uniform tiles. When a later tile has content already in the map,
\textbf{do not write its body again}: instead set its entry in the format's
tile-offset and tile-bytecount arrays to point at the \emph{same} byte range as
the earlier tile.

The result is that each pyramid level of the file stores exactly one
compressed body for \enquote{all-zero}, one for \enquote{all-one}, and so on,
and the hundreds of thousands of ocean tiles at that level all reference it.
It is unusual for multiple image tiles to resolve to one shared byte range ---
most encoders never do it --- but the TIFF~6.0 specification permits it (the
\code{TileOffsets} and \code{TileByteCounts} arrays are just arrays of pointers
and lengths; nothing requires them to be distinct or monotone), and mainstream
readers (e.g.\ GDAL, \code{tifffile}, browser GeoTIFF libraries) decode it
correctly.

A corollary: the file's structural metadata must \textbf{not} advertise a
strict tile ordering or contiguous tile layout (for COGs, the \code{BLOCK\_ORDER}
ghost option is deliberately omitted), because such a declaration would forbid
tiles from sharing bytes.

\emph{Reader-side complement.} A client that decodes tiles on demand should key
its decoded-tile cache by each tile's byte offset in the container rather than
by grid position. Because the technique above makes many grid positions resolve
to the same byte range, an offset-keyed cache collapses every reference to a
shared body --- for example all of the ocean tiles at a given level --- into a
single entry, so scanning a large uniform region costs one cache slot instead
of repeatedly evicting useful tiles. The client libraries of \S6 do this.

\subsubsection{Monotone value transform}

Cell values are stored as a fixed monotone transform of the density, for
example $\ln(1 + \text{density})$. This preserves the ordering of cells (so any
threshold or ranking is unchanged) while compressing a very skewed dynamic
range (roughly $0$ to $2.3 \times 10^{5}$) into a near-linear small range
(roughly $0$ to $12$), which improves both the compressibility of tile bodies
and the out-of-the-box behavior of visualization tools. The transform's
maximum over the raster is recorded in a standard tag so that readers can
invert or re-normalize it.

\subsection{Cloud-Optimized GeoTIFF assembly}

A COG places all the image file directories (IFDs) and a small block of
structural \enquote{ghost} metadata ahead of the tile data, so that a reader
learns the file's structure from a short prefix and then fetches only the tiles
it needs with HTTP range requests. But the tile offsets are not known until the
tiles have been written. The method therefore:

\begin{enumerate}
\item streams compressed tile bodies to a temporary file during the pass
  (\S3.4), recording each body's provisional offset within that file and its
  length;
\item writes the file header, the ghost metadata, and all IFDs (full-resolution
  level and every overview level), with placeholder tile-offset values;
\item for each pyramid level in turn, appends that level's tile-offset array
  followed by its tile bodies, drawn from the temporary file and now at their
  final absolute offsets; the per-level tile-bytecount arrays follow at the end
  of the file;
\item seeks back and patches every IFD's tile-offset and tile-bytecount
  pointer, and the inter-IFD \enquote{next IFD} pointers, to the final values.
\end{enumerate}

Classic (32-bit) TIFF is sufficient: the size reductions of \S3.6 keep the
output below 4~GiB. Standard geo-referencing tags place the grid in the target
projection; provenance tags record the input date range, the producing software
version, and the date of the most recent input period.

%========================================================================
\section{Resulting properties}
%========================================================================

\begin{itemize}
\item \textbf{Peak memory constant in cell count.} Determined by tree depth
  $L$, window size $N$ (one decompressing reader per period), and compression
  buffers. A worked configuration builds a global zoom-18 raster with pyramid
  in about 800~MiB resident, well under 1~GiB.
\item \textbf{Single pass.} The input is read once; overviews are a by-product;
  no stage re-reads the full-resolution raster.
\item \textbf{Output size tracks information, not dimensions.} A worked
  configuration emits ${\approx}\,600$~MB --- the exact figure drifts from week
  to week with the input --- where a dense zoom-18 encoding would be
  ${\approx}\,275$~GB.
\item \textbf{Incremental.} Re-running for a shifted window re-sorts only the
  new period(s); everything else is reused from cache.
\item \textbf{Standards-compliant output.} A valid Cloud-Optimized GeoTIFF
  readable by unmodified mainstream tools, including the shared-offset tiles of
  \S3.6.2.
\item \textbf{Commodity hardware.} No cluster, no large-memory machine, no GPU.
\end{itemize}

%========================================================================
\section{Variations and generalizations}
%========================================================================

The disclosure covers, without limitation, the following variations:

\begin{itemize}
\item \textbf{Any raster produced by aggregating a key-ordered time series into
  a tiled pyramid}, not only view-density maps but also population, emissions,
  elevation change, sensor coverage, edit activity, and the like.
\item \textbf{Any spatial tree.} Quad-trees (this document), binary tile
  pyramids, octrees / 3D tiles for volumetric or point data, or non-square
  tilings, by adjusting the bits-per-level of the key of \S3.1 and the arity of
  the eviction/subsample step.
\item \textbf{Any projection or CRS}, geographic or projected; the
  area-normalization of \S3.3 is projection-specific but optional.
\item \textbf{Any robust per-cell reduction} computed from an ascending run
  plus an absent-period count: other quantiles, trimmed or winsorized means,
  decay- or recency-weighted means, min, max.
\item \textbf{Any pooling rule} for overview generation: max, min, mean,
  median, mode, or a learned kernel.
\item \textbf{Any container format} whose tile index is an array of
  (offset,~length) pairs and does not mandate distinctness --- TIFF, BigTIFF,
  GeoTIFF, or a bespoke format --- for the deduplication of \S3.6.2; use BigTIFF
  in place of classic TIFF when \S3.6 does not keep the file under 4~GiB.
\item \textbf{Any lossless codec} for tile bodies; \textbf{any monotone
  transform} (logarithm, square root, $\mu$-law, a fitted CDF) for \S3.6.3.
\item \textbf{Any spill-to-disk external sort and any $k$-way merge} for
  \S3.2, including merging directly from remote object storage.
\item \textbf{Gap fill} by zero, by the nearest painted ancestor, by
  interpolation, or by leaving an explicit nodata value.
\item \textbf{Parallelization.} The per-period sorts of \S3.2 are trivially
  parallelizable. The merge-and-paint pass of \S3.3--\S3.7 can be sharded by
  sub-tree: each shard covers a disjoint key range and quad-tree sub-tree,
  painting it and its overviews down to the shard root; a short serial tail
  then subsamples the shard roots up the remaining pyramid levels, concatenates
  the tile bodies (whose order is immaterial), and fixes up the IFD offsets.
  Because sub-trees over ocean carry almost no data while urban ones carry most
  of it, shard far more finely than the core count and balance the shards
  dynamically. Within a shard, only the pre-order walk and its
  eviction/subsample step are inherently sequential; input decompression (one
  worker per period file) and tile-body compression (one per evicted tile)
  parallelize freely.
\item \textbf{Deferred / range-served overviews.} Emitting fewer overview
  levels, or none, and relying on the reader.
\item \textbf{Reader-side caching by shared offset.} A client that reads a
  container with deduplicated tiles (\S3.6.2) keys its decoded-tile cache on
  each tile's file offset rather than its grid coordinates, so every reference
  to one shared body occupies a single cache entry.
\item \textbf{Bounded reads from fixed tile geometry.} Because every tile
  decodes to exactly $T \times T \times \operatorname{sizeof}(\text{cell})$
  bytes, a reader can cap the decompressor's output at that size and reject,
  when opening the file, any tile whose recorded compressed length exceeds a
  small margin over the largest plausible compressed tile --- bounding memory
  use on adversarial input with no separate size field.
\end{itemize}

%========================================================================
\section{Reference implementation}
%========================================================================

A complete, working implementation in the Go programming language is publicly
available under the MIT license at \url{https://github.com/brawer/osmviews}
(directory \code{cmd/osmviews-builder}; the tile key is in \code{tile.go}, the
two-stage ordering in \code{tilelogs.go} and \code{merge.go}, the bounded
raster path and inline overviews in \code{paint.go}, and the size reductions
and COG assembly in \code{raster.go}).
Reader-side reference implementations --- client libraries that query the
resulting file and implement the offset-keyed cache of \S3.6.2 --- are
published under the MIT license at
\url{https://github.com/brawer/osmviews-rs} (Rust) and
\url{https://github.com/brawer/osmviews-py} (Python). The commit history of
these repositories predates or accompanies this publication.

%========================================================================
\phantomsection
\addcontentsline{toc}{section}{Appendix: enumerated disclosed techniques}
\section*{Appendix: enumerated disclosed techniques}
%========================================================================

For the avoidance of doubt, the following are disclosed for free public use,
alone and in combination:

\begin{enumerate}
\item Encoding mixed-level map tiles as integers whose ascending numeric order
  is a pre-order traversal of the tiling tree, by reserving low (or high) bits
  for the level and interleaving coordinate bits per level.
\item Aggregating a windowed time series of per-tile counts by (a)~an external,
  disk-backed sort of each period on that key, then (b)~a streaming $k$-way
  merge across periods, yielding one key-ordered stream in which all periods'
  records for a tile are consecutive.
\item Applying a per-period scalar weight to counts during the merge to
  normalize periods with missing sub-intervals.
\item Computing a per-cell windowed order statistic (e.g.\ median with
  zero-fill for absent periods) directly from the ascending per-cell run and
  the count of absent periods, without materializing the window for all cells.
\item Painting the aggregate into a tree of fixed-size raster tiles while
  holding only the root-to-current-leaf path, justified by the pre-order
  arrival of records, with immediate eviction of finalized raster tiles.
\item Generating the entire reduced-resolution overview pyramid inline, by
  subsampling each raster tile into its parent at eviction time, in the same
  single pass over the input.
\item Emitting skipped sub-trees and never-touched regions as uniform raster
  tiles so the output has global coverage with no holes.
\item Quantizing raster tiles to a common value at storage precision so that
  large low-magnitude regions collapse to identical tiles.
\item Deduplicating identical raster tiles by setting multiple entries of the
  container's tile-offset / tile-length index to point at a single shared byte
  range, and correspondingly omitting any strict-tile-order declaration from
  the container's structural metadata.
\item Storing a monotone transform of the cell value to compress dynamic range
  while preserving order, with the transform's extent recorded in metadata.
\item Assembling a header-first (Cloud-Optimized) GeoTIFF by staging tile
  bodies in a temporary file and back-patching offsets after the single pass.
\item The combination of 1--11 to build a globally complete, high-zoom,
  overview-complete Cloud-Optimized GeoTIFF from tile-access logs using working
  memory that is constant in the number of cells.
\item Reading such a container by caching decoded tiles keyed on their shared
  byte offset, so that all references to a deduplicated tile body occupy one
  cache entry.
\item Bounding a reader's memory use by capping decompression output at the
  fixed decoded tile size and rejecting, at open time, any tile whose recorded
  compressed length exceeds the largest plausible compressed tile.
\item The combination of 1--14 as an end-to-end scheme for producing such a
  container and reading it.
\end{enumerate}

\end{document}
